\subsection{An explicit sensitivity matrix for every stored state}
\label{app:explicit-state-matrix}
The following construction replaces unspecified composition constants by
finite algebraic expressions. Norms compare stored QR frames, not merely
projectors. We specify one routed matrix; further routed layers or fallback
states require additional rows and columns. Moment states start at zero,
so fresh bias-corrected first moments are bounded by $G_0$ as in
Lemma~\ref{lem:actual-invariants}. Let $e_W,e_U,e_M,e_A,e_H,e_g$ be the coordinate row vectors of
$\mathbb R^6$ and
\begin{equation}
 d=(\delta W,\delta U,\delta M,\delta A,\delta H,\delta g)^{\top}\ge0,
 \qquad\|G-G'\|\le L_g\delta W+\delta g.
\end{equation}
For a common example $\delta g=0$; for one replaced example at a common
state $\delta g\le2G_0$. Assume all reachable gradients are bounded by
$G_0$, all raw first moments and root second moments by $G_0$, and all
tracker momenta by $H_0$. The QR comparison segments have minimum singular
value at least $s_Q>0$; Procrustes inputs have minimum singular value at
least $s_R>0$. Define $\kappa_Q=\sqrt3/s_Q$,
$\kappa_R=1/s_R$, and the following nonnegative rows:
\begin{align}
 k_G&=L_ge_W+e_g,\\
 k_Z&=k_G+G_0e_U, & k_E&=k_G+2G_0e_U,\\
 k_M&=\beta_1e_M+(1-\beta_1)k_Z,\\
 k_A&=\sqrt{\beta_2}e_A+\sqrt{1-\beta_2}k_Z,\\
 k_D&=\frac{\varepsilon_a^{-1}}{1-\beta_1^t}k_M
       +\frac{G_0\varepsilon_a^{-2}}{\sqrt{1-\beta_2^t}}k_A
       +G_0\varepsilon_a^{-1}e_U,\\
 k_{\rm res}&=(s_{\max}+G_0^2\varepsilon_r^{-2})k_E
                          +G_0\varepsilon_r^{-1}k_G,\\
 k_Q&=4e_U+2\varepsilon_c^{-1/2}k_G,\\
 k_H&=\mu_He_H+2\mu_HH_0e_U+(1-\mu_H)k_Q,\\
 k_{U+}&=\kappa_Q(e_U+\rho_t k_H),\\
 k_R&=\kappa_R(k_{U+}+e_U),\\
 k_{M+}&=k_M+G_0k_R,& k_{A+}&=k_A+G_0k_R,\\
 k_{W+}&=|1-\eta_t\lambda_{\rm wd}|e_W
                         +\eta_t k_D+\eta_{\perp,t} k_{\rm res}.
 \label{eq:all-state-rows}
\end{align}
Here $k_{W+},k_{U+},k_{M+},k_{A+},k_H$ are stacked into
$\mathbb K_t\in\mathbb R_+^{5\times6}$; write
$\mathbb K_t=(K_t\mid v_t)$.

\begin{proposition}[Constructive full-state Lipschitz and injection constants]
\label{prop:explicit-state-constants}
Under the preceding bounds, the one-step state differences obey
\begin{equation}
 \begin{pmatrix}\delta W_+\\\delta U_+\\\delta M_+\\\delta A_+\\\delta H_+\end{pmatrix}
 \le K_t\begin{pmatrix}\delta W\\\delta U\\\delta M\\\delta A\\\delta H\end{pmatrix}
                       +v_t\delta g.\label{eq:matrix-state-bound}
\end{equation}
For $\omega=(1,\omega_U,\omega_M,\omega_A,\omega_H)>0$,
valid constants in Assumption~\ref{ass:stability} are
\begin{equation}
 L_t=\max_j\frac{\sum_i\omega_i(K_t)_{ij}}{\omega_j},\qquad
 b_t=2G_0\sum_i\omega_i(v_t)_i.\label{eq:explicit-L-b}
\end{equation}
In particular, identical input states and identical examples imply zero
output difference exactly.
\end{proposition}
\begin{proof}
The routed and residual gradient estimates are
\begin{align}
 \|GU-G'U'\|&\le\|G-G'\|+G_0\|U-U'\|,\\
 \|G(I-P)-G'(I-P')\|
 &\le\|G-G'\|+2G_0\|U-U'\|.
\end{align}
The first-moment EMA gives $k_M$. Reverse triangle inequalities for the
coordinatewise two-dimensional Euclidean norms give $k_A$. Bias correction
and \eqref{eq:adam-lipschitz}, followed by lifting, give $k_D$.
For $n=\|G\|$, $e=\|G(I-P)\|$, scalar clipping is nonexpansive, so
\begin{align}
 |s-s'|&\le\varepsilon_r^{-1}|n-n'|
                        +G_0\varepsilon_r^{-2}|e-e'|,\\
 \|sE-s'E'\|&\le s_{\max}\|E-E'\|+G_0|s-s'|.
\end{align}
These are precisely the coefficients in $k_{\rm res}$.
Lemma~\ref{lem:building-blocks} gives $k_Q$, while
\begin{equation}
 \|(I-P)H-(I-P')H'\|\le\|H-H'\|+2H_0\|U-U'\|
\end{equation}
gives $k_H$. Lemma~\ref{lem:qr-frame} gives $k_{U+}$.
The polar perturbation inequality gives
\begin{equation}
 \|R-R'\|\le s_R^{-1}
 \|U_+^\top U-U_+'{}^\top U'\|
 \le\kappa_R(\delta U_++\delta U).
\end{equation}
Finally,
\begin{align}
 \|\widetilde MR^\top-\widetilde M'R'^\top\|
 &\le\|\widetilde M-\widetilde M'\|+G_0\|R-R'\|,\\
 \|\sqrt{\widetilde A^{\odot2}(R^\top)^{\odot2}}
       -\sqrt{\widetilde A'{}^{\odot2}(R'^\top)^{\odot2}}\|
 &\le\|\widetilde A-\widetilde A'\|+G_0\|R-R'\|.
\end{align}
The second inequality follows by writing each output amplitude as the norm
of a row weighted by a column of the rotation, then using its orthogonality.
These prove every row of \eqref{eq:all-state-rows}. Multiplying
\eqref{eq:matrix-state-bound} by $\omega^\top$ yields
\eqref{eq:explicit-L-b}. All operations are deterministic functions of
states and examples; there is no sample-independent additive error.
\end{proof}

\begin{corollary}[Matrix-valued replacement stability]
For independent uniform sampling with replacement and common data-independent
initialization, define
$u_t=\mathbb E(\delta W_t,\delta U_t,\delta M_t,\delta A_t,\delta H_t)^\top$.
If the constants above are uniform over the two coupled runs, then
\begin{align}
 u_{t+1}&\le K_tu_t+\frac{2G_0}{n}v_t,\\
 u_T&\le\frac{2G_0}{n}\sum_{t=0}^{T-1}
                  K_{T-1}\cdots K_{t+1}v_t,\label{eq:matrix-stability-product}\\
 |\mathbb E[L(W_T)-\widehat L_S(W_T)]|
 &\le\frac{2L_\ell G_0}{n}
 e_W^{(5)}\sum_{t=0}^{T-1}K_{T-1}\cdots K_{t+1}v_t.
\end{align}
Here $e_W^{(5)}=(1,0,0,0,0)$ and an empty product is the identity.
\end{corollary}
\begin{proof}
The probability of drawing the replaced example is $1/n$; iterate
\eqref{eq:matrix-state-bound} from $u_0=0$ and use the ghost-sample
identity in Theorem~\ref{thm:generalization}.
\end{proof}

\subsection{Spectral sensitivity, feedback, and a verifiable contractive specialization}
\label{app:spectral-stability}

\begin{proposition}[Optimal weights and the Schur feedback condition]
\label{prop:spectral-feedback}
For a nonnegative time-independent sensitivity matrix $K$,
\begin{equation}
 \inf_{\omega>0}\|D_\omega KD_\omega^{-1}\|_1
 =\operatorname{spr}(K).\label{eq:perron-weights}
\end{equation}
Partition a valid sensitivity bound as
\begin{equation}
 K=\begin{pmatrix}k&b^\top\\c&A\end{pmatrix},\qquad
 k\ge0,\quad A,b,c\ge0,\quad\operatorname{spr}(A)<1.
\end{equation}
Then
\begin{equation}
 \operatorname{spr}(K)<1
 \quad\Longleftrightarrow\quad
 k+b^\top(I-A)^{-1}c<1.\label{eq:schur-feedback}
\end{equation}
For an injection $h=(h_W,h_S)^\top\ge0$ and
$u_{t+1}\le Ku_t+h/n$, $u_0=0$, the resulting uniform parameter sensitivity
satisfies
\begin{equation}
 \sup_T(u_T)_W
 \le\frac{h_W+b^\top(I-A)^{-1}h_S}
 {n\{1-k-b^\top(I-A)^{-1}c\}}.\label{eq:schur-stability}
\end{equation}
The triangle-inequality construction \eqref{eq:all-state-rows} usually has
$k\ge1$ when weight decay is absent. It does not by itself establish
\eqref{eq:schur-feedback}; that requires a sharper signed contraction
estimate for the parameter transition.
\end{proposition}
\begin{proof}
Every induced norm bounds the spectral radius from below. For any
$q>\operatorname{spr}(K)$,
\begin{align}
 \omega&=(qI-K^\top)^{-1}\mathbf1
       =\sum_{j=0}^{\infty}q^{-j-1}(K^\top)^j\mathbf1>0,\\
 K^\top\omega&=q\omega-\mathbf1<q\omega.
\end{align}
Hence the norm in \eqref{eq:perron-weights} is less than $q$;
let $q$ decrease to $\operatorname{spr}(K)$.
If \eqref{eq:schur-feedback} holds, choose
\begin{equation}
 x=\binom{1}{(I-A)^{-1}c+\epsilon(I-A)^{-1}\mathbf1},\quad\epsilon>0.
\end{equation}
For sufficiently small $\epsilon$, $Kx<x$, proving
$\operatorname{spr}(K)<1$. Conversely, if $\operatorname{spr}(K)<1$,
then $(I-K)^{-1}=\sum_{j\ge0}K^j\ge0$. Its top-left entry, computed by
Schur complementation, is
$[1-k-b^\top(I-A)^{-1}c]^{-1}$ and is positive, proving the converse.
Finally $u_T\le n^{-1}\sum_{j=0}^{T-1}K^jh\le n^{-1}(I-K)^{-1}h$;
the top component is \eqref{eq:schur-stability}.
\end{proof}

\begin{proposition}[A signed, frozen-metric quadratic specialization]
\label{prop:quadratic-stability}
Let
$\ell(w,z)=w^\top H w/2-b_z^\top w$, $H\succ0$, $\|b_z\|\le B_z$,
and use a fixed positive metric $B\succ0$ and momentum $\beta\in[0,1)$:
\begin{equation}
 m_t=\beta m_{t-1}+(1-\beta)(Hw_t-b_{I_t}),\qquad
 w_{t+1}=w_t-\eta Bm_t.
\end{equation}
Set $\widetilde H=B^{1/2}HB^{1/2}$, with eigenvalues
$0<\gamma_{\min}\le\gamma_j\le\gamma_{\max}$.
The signed two-state transition is asymptotically stable if
\begin{equation}
 0<\eta<\frac{2(1+\beta)}{(1-\beta)\gamma_{\max}}.\label{eq:quadratic-schur}
\end{equation}
For any $q\in(\max_j\operatorname{spr}(T_{\gamma_j}),1)$ there is a
finite $C_q$ such that the neighboring-data parameter difference obeys
\begin{equation}
 \sup_T\mathbb E\|w_T-w_T'\|
 \le\frac{2C_q(1-\beta)B_z\|B^{1/2}\|_{\rm op}^2
                           \sqrt{1+\eta^2}}
                         {n(1-q)}.\label{eq:quadratic-generalization}
\end{equation}
This is a frozen-metric specialization, not a contraction theorem for
unrestricted Adam or a moving routing basis.
\end{proposition}
\begin{proof}
Use coordinates $x=B^{-1/2}w$, $y=B^{1/2}m$ and diagonalize
$\widetilde H$. Each eigencoordinate has transition
\begin{equation}
 T_\gamma=\begin{pmatrix}
 1-\eta(1-\beta)\gamma&-\eta\beta\\
 (1-\beta)\gamma&\beta
 \end{pmatrix},\qquad
 \det(zI-T_\gamma)=z^2-[1+\beta-\eta(1-\beta)\gamma]z+\beta.
\end{equation}
The real second-order Schur conditions reduce to
\begin{align}
 1-\beta&>0,\\
 1-[1+\beta-\eta(1-\beta)\gamma]+\beta
 &=\eta(1-\beta)\gamma>0,\\
 1+[1+\beta-\eta(1-\beta)\gamma]+\beta
 &=2(1+\beta)-\eta(1-\beta)\gamma>0.
\end{align}
These prove \eqref{eq:quadratic-schur}. In finite dimension,
$\|T^k\|\le C_qq^k$ for any $q$ strictly larger than its spectral radius,
including nondiagonalizable blocks. A replaced example injects a vector
of norm at most
$2(1-\beta)B_z\|B^{1/2}\|_{\rm op}\sqrt{1+\eta^2}$ into $(x,y)$
with probability $1/n$. Sum the geometric response and multiply by
$\|B^{1/2}\|_{\rm op}$ to return to $w$.
\end{proof}

\subsection{Validation selection is a separate statistical operation}
\label{app:validation-selection}
\begin{theorem}[Finite-candidate validation selection]
Let $\theta_1(S),\ldots,\theta_K(S)$ be any predictors constructed without
access to an independent validation set $V$ of size $n_v$. For losses in
$[0,1]$, let $\widehat k\in\arg\min_k\widehat L_V(\theta_k)$.
With probability at least $1-\delta$ over $V$, conditional on $S$ and all
candidate-generation randomness,
\begin{align}
 \max_{k\le K}|L(\theta_k)-\widehat L_V(\theta_k)|
 &\le\epsilon_v:=\sqrt{\frac{\log(2K/\delta)}{2n_v}},\\
 L(\theta_{\widehat k})&\le\min_{k\le K}L(\theta_k)+2\epsilon_v.
 \label{eq:validation-oracle}
\end{align}
Every inspected checkpoint, learning rate, rank, and restart must be
included among the candidates. Repeated adaptive use of $V$ to generate
new candidates requires a different argument.
\end{theorem}
\begin{proof}
Conditional on the candidates, Hoeffding and the union bound yield
\begin{equation}
 \Pr\!\left(\max_k|L_k-\widehat L_{V,k}|>u\mid S,\Omega\right)
 \le2K e^{-2n_vu^2}.
\end{equation}
On its complementary event, for every $k$,
\begin{equation}
 L_{\widehat k}\le\widehat L_{V,\widehat k}+\epsilon_v
 \le\widehat L_{V,k}+\epsilon_v\le L_k+2\epsilon_v.
\end{equation}
This proof neither assumes nor transfers training-set stability of the
selected learning-rate map.
\end{proof}

\subsection{PAC-Bayes: anisotropic posterior optimization and a margin specialization}
\label{app:pacbayes-derivations}
Let $\Pi$ be data independent and $\ell\in[0,1]$. For fixed $\lambda>0$,
Hoeffding's lemma and independence imply
\begin{align}
 \mathbb E_S\exp\{\lambda[L(\theta)-\widehat L_S(\theta)]\}
 &\le e^{\lambda^2/(8n)},\\
 \mathbb E_S\mathbb E_{\theta\sim\Pi}
 \exp\{\lambda[L(\theta)-\widehat L_S(\theta)]-\lambda^2/(8n)\}
 &\le1.\label{eq:pb-mgf-explicit}
\end{align}
Markov's inequality and the Gibbs--Donsker--Varadhan variational principle
\cite{donsker1975asymptotic} give,
on a single event of probability at least $1-\delta$, simultaneously for
all $Q\ll\Pi$,
\begin{equation}
 L(Q)\le\widehat L_S(Q)+
 \frac{\operatorname{KL}(Q\|\Pi)+\log(1/\delta)}{\lambda}
                     +\frac{\lambda}{8n}.\label{eq:pb-uniform-Q}
\end{equation}
Indeed, for $F=\lambda(L-\widehat L_S)$,
\begin{align}
 \mathbb E_QF-\operatorname{KL}(Q\|\Pi)
 &=\mathbb E_Q\log\left(e^F\frac{d\Pi}{dQ}\right)
 \le\log\mathbb E_\Pi e^F.
\end{align}
For a countable data-independent set $\{\lambda_j\}$ with masses
$\pi_j>0$, $\sum_j\pi_j=1$, replace $\log(1/\delta)$ by
$\log[1/(\delta\pi_j)]$ to obtain simultaneous validity and permit
selection of $j$ after training.

\begin{proposition}[Optimal two-variance Gaussian certificate]
\label{prop:optimal-posterior}
Let $P=UU^\top$ be any data-dependent rank-$r$ projector, $0<r<d$, and
$\Pi=\mathcal N(\operatorname{vec}W_0,s_0^2I_{md})$, with $W_0,s_0$
data independent. Let
\begin{equation}
 Q=\mathcal N(\operatorname{vec}W,
          I_m\otimes[s_\parallel^2P+s_\perp^2(I-P)]),
 \quad s_\parallel,s_\perp>0,
\end{equation}
using row-wise vectorization. Suppose for all $E$ and $u\in[0,1]$,
\begin{equation}
 \langle E,\nabla^2\widehat L_S(W+uE)[E]\rangle
 \le h_\parallel\|EP\|^2+h_\perp\|E(I-P)\|^2,
 \quad h_\parallel,h_\perp\ge0.\label{eq:posterior-hessian}
\end{equation}
Then the variances minimizing the sum of the quadratic empirical-loss
upper bound and $\operatorname{KL}/\lambda$ are
\begin{equation}
 s_\parallel^2=(s_0^{-2}+\lambda h_\parallel)^{-1},\qquad
 s_\perp^2=(s_0^{-2}+\lambda h_\perp)^{-1}.\label{eq:optimal-variances}
\end{equation}
On the event \eqref{eq:pb-uniform-Q}, this posterior satisfies
\begin{align}
 L(Q)\le\widehat L_S(W)
 &+\frac{\|W-W_0\|^2}{2\lambda s_0^2}
 +\frac{m}{2\lambda}\{r\log(1+\lambda h_\parallel s_0^2)
                 +(d-r)\log(1+\lambda h_\perp s_0^2)\}\notag\\
 &+\frac{\log(1/\delta)}{\lambda}+\frac{\lambda}{8n}.
 \label{eq:optimized-pb}
\end{align}
No derivative of the learned projector with respect to the sample is
needed for this PAC-Bayes statement; the prior must remain data independent.
\end{proposition}
\begin{proof}
Put $\Sigma=s_\parallel^2P+s_\perp^2(I-P)$. Its eigenvalues and determinant
are
\begin{align}
 \operatorname{spec}\Sigma
 &=\{s_\parallel^2\ (r\text{ times}),s_\perp^2\ (d-r\text{ times})\},\\
 \operatorname{tr}\Sigma&=rs_\parallel^2+(d-r)s_\perp^2,\\
 \log\det\Sigma&=r\log s_\parallel^2+(d-r)\log s_\perp^2.
\end{align}
The Gaussian density ratio integrates to
\begin{align}
 2\operatorname{KL}(Q\|\Pi)
 &=s_0^{-2}\|W-W_0\|^2
   +m\{s_0^{-2}\operatorname{tr}\Sigma-d
                         +d\log s_0^2-\log\det\Sigma\}\\
 &=s_0^{-2}\|W-W_0\|^2
   +mr\,k(s_\parallel^2/s_0^2)
   +m(d-r)k(s_\perp^2/s_0^2),\quad k(x)=x-1-\log x.
\end{align}
Taylor's integral formula and $\mathbb E E=0$ imply
\begin{align}
 \widehat L_S(Q)
 &\le\widehat L_S(W)+\tfrac12m
                    [rh_\parallel s_\parallel^2+(d-r)h_\perp s_\perp^2].
\end{align}
For a subspace of multiplicity $p$, minimize over $v>0$:
\begin{align}
 \phi_h(v)&=\tfrac{mp}{2}\{hv+\lambda^{-1}k(v/s_0^2)\},\\
 \phi_h'(v)&=\tfrac{mp}{2}\{h+\lambda^{-1}(s_0^{-2}-v^{-1})\},\\
 \phi_h''(v)&=\frac{mp}{2\lambda v^2}>0,\\
 \min_{v>0}\phi_h(v)&=\frac{mp}{2\lambda}\log(1+\lambda hs_0^2).
\end{align}
Apply this separately to $p=r,d-r$ and substitute into
\eqref{eq:pb-uniform-Q}. The Hessian condition is global along Gaussian
perturbation segments; a local empirical Hessian estimate alone does not
verify it for an untruncated Gaussian posterior.
\end{proof}

\begin{corollary}[A deterministic center under a loss Lipschitz condition]
If $\ell(\cdot,z)$ is $L_\ell$-Lipschitz for every $z$, then
\begin{equation}
 L(W)\le L(Q)+L_\ell
         \sqrt{m[rs_\parallel^2+(d-r)s_\perp^2]}.\label{eq:pb-center}
\end{equation}
\end{corollary}
\begin{proof}
$L(W)-L(Q)\le L_\ell\mathbb E_Q\|W-\widetilde W\|
\le L_\ell\sqrt{\mathbb E_Q\|W-\widetilde W\|^2}$.
\end{proof}

\begin{proposition}[Gaussian conditioning and a sharp center-to-Gibbs conversion]
\label{prop:pb-margin}
Let $x=\phi(z)\ne0$ be a feature map fixed before observing $S$, and
let $W\in\mathbb R^{m\times d}$, $m\ge2$, be a multiclass linear head.
Use the row-independent Gaussian posterior
$Q=\mathcal N(\operatorname{vec}W,I_m\otimes\Sigma)$, $\Sigma\succ0$,
and a legal data-independent prior $\Pi$. Write $\varphi,\Phi$ for the
standard normal density and distribution function. For each training pair,
\begin{align}
 \nu_i^2&=x_i^\top\Sigma x_i,&
 \Delta_{ij}&=(Wx_i)_{y_i}-(Wx_i)_j,\\
 p_i&=\int_{\mathbb R}\varphi(t)
       \left[1-\prod_{j\ne y_i}\Phi(t+\Delta_{ij}/\nu_i)\right]dt,&
 p_Q&=n^{-1}\sum_i p_i.
\end{align}
The empirical Gibbs classification error is exactly $p_Q$. With
probability at least $1-\delta$,
\begin{equation}
 L_{0/1}(W)\le\min\left\{1,
 2\operatorname{kl}^{-1}_+\left(p_Q,
      \frac{\operatorname{KL}(Q\|\Pi)+\log(\Xi_n/\delta)}n\right)\right\}.
 \label{eq:pb-margin}
\end{equation}
For a simpler empirical margin summary, define
\begin{align}
 \widehat L_\gamma(W)&=n^{-1}\sum_i
                  \mathbf1\{\min_{j\ne y_i}\Delta_{ij}\le\gamma\},\\
 B_m(h)&=\int_{\mathbb R}\varphi(t)[1-\Phi(t+h)^{m-1}]\,dt,
 \qquad\nu_{\max}=\max_{i\le n}\nu_i.
\end{align}
For $\gamma>0$, $p_Q$ in \eqref{eq:pb-margin} may be replaced by
\begin{equation}
 \overline p_\gamma=\widehat L_\gamma(W)
       +(1-\widehat L_\gamma(W))B_m(\gamma/\nu_{\max}).
\end{equation}
In particular $B_2(h)=1-\Phi(h/\sqrt2)$, and
$B_m(h)\le\min\{1,(m-1)[1-\Phi(h/\sqrt2)]\}$ for $h\ge0$.
No global feature-norm bound is needed for this empirical summary.
\end{proposition}
\begin{proof}
Conditional on the true-class perturbation $E_{y_i}x_i=\nu_it$, the other
logit perturbations are independent $\mathcal N(0,\nu_i^2)$. The conditional
probability of correct classification is
$\prod_{j\ne y_i}\Phi(t+\Delta_{ij}/\nu_i)$; integrate to get $p_i$.
If all mean margins exceed $\gamma$, monotonicity gives
$p_i\le B_m(\gamma/\nu_i)\le B_m(\gamma/\nu_{\max})$.
For bad-margin examples use $p_i\le1$, retaining the exact fraction of good
examples. The binary and pairwise bounds use
$E_jx-E_yx\sim\mathcal N(0,2\nu(x)^2)$ and the union bound; no
coordinatewise-infinity-norm relaxation is needed.

For every test input misclassified by $W$, select a competitor with mean
logit at least the true-class logit. Its centered symmetric Gaussian logit
difference makes a Gibbs error with probability at least $1/2$. Hence
$L_{0/1}(W)\le2L_{0/1}(Q)$. This factor two is attained by a binary
zero-mean-margin example when ties are counted as center errors.
Theorem~\ref{thm:Seeger-Maurer} and monotonicity of the upper KL inverse
prove \eqref{eq:pb-margin}. The identities use an independent fixed feature
map; an encoder learned from $S$ must be included in the hypothesis space
and prior construction. Numerical evaluation of the integrals requires a
controlled quadrature error if used for a certificate.
\end{proof}
